% Chapter 2: Related Works
\subsection{Information Theory Primer}
% • Entropy coding and basic formulations
% • Distributed coding: Slepian–Wolf (lossless), Wyner–Ziv (lossy)
% • Conventional vs. learned DSC methods

Information theory, pioneered by Claude Shannon in the mid-20th century \cite{6773024}, provides a mathematical framework for quantifying information, communication, and data compression. At its core, the field seeks to answer two fundamental questions: What is the ultimate limit of data compression (source coding), and what is the ultimate limit of the rate of reliable communication over a noisy channel (channel coding)? This primer will focus on the principles of source coding in a distributed context, which are most relevant to this thesis.

\subsubsection{Entropy and Mutual Information}
At the core of information theory is the concept of \emph{entropy}, which measures the uncertainty or average information content in a random variable. For a discrete random variable $X$ with possible outcomes $x_1, x_2, \dots, x_n$ and probability mass function $P(X = x_i) = p_i$, the Shannon entropy $H(X)$ is defined as:
\begin{equation}
    H(X) = -\sum_{i=1}^n p_i \log_2 p_i,
\end{equation}
where the logarithm is base-2, yielding entropy in units of bits per symbol. Entropy represents the minimum average number of bits required to encode outcomes of $X$ without loss of information. For instance, a fair coin flip has an entropy of 1 bit, representing the maximum uncertainty for a binary event. A biased coin, which lands on heads 90\% of the time, has a much lower entropy because the outcome is more predictable. For continuous random variables, the \emph{differential entropy} $h(X)$ is analogous:

\begin{equation}
    h(X) = -\int_{-\infty}^{\infty} f(x) \log_2 f(x) \, dx,
\end{equation}
where $f(x)$ is the probability density function. However, because of the continuous nature of the variable, differential entropy can be negative and is not directly comparable to discrete entropy. Furthermore, for two random variables $X$ and $Y$, the \emph{join entropy} $H(X,Y)$ is defined as:
\begin{equation}
    H(X,Y) = -\sum_{x,y} p(x,y) \log_2 p(x,y),
\end{equation}
which quantifies the total uncertainty in the pair $(X,Y)$. It satisfies the logical condition $H(X,Y) \geq \max(H(X), H(Y))$, with equality if one variable is a deterministic function of the other. Similarly, \emph{Conditional entropy} $H(X|Y)$ measures the remaining uncertainty in $X$ given knowledge of $Y$:
\begin{equation}
    H(X|Y) = H(X,Y) - H(Y) = -\sum_{x,y} p(x,y) \log_2 p(x|y) = \sum_y p(y) H(X|Y=y).
\end{equation}

Intuitively, $H(X|Y)$ represents the average additional bits needed to describe $X$ when $Y$ is already known. It is always non-negative, $H(X|Y) \geq 0$, and satisfies $H(X|Y) \leq H(X)$, with equality if $X$ and $Y$ are independent (i.e., knowing $Y$ provides no reduction in uncertainty about $X$). If $X$ is a deterministic function of $Y$, then $H(X|Y) = 0$. The chain rule for entropy extends this: $H(X,Y) = H(Y) + H(X|Y) = H(X) + H(Y|X)$, allowing decomposition of joint entropies for multiple variables. For $n$ variables:
\begin{equation}
    H(X_1, \dots, X_n) = H(X_1) + \sum_{i=2}^n H(X_i | X_1, \dots, X_{i-1})
\end{equation}
which decomposes joint entropy into a sum of conditional entropies. This is particularly useful in analyzing sequential or dependent processes, such as Markov chains. %, where $H(X_n|X_1, \dots, X_{n-1}) = H(X_n|X_{n-1})$ under the Markov property.
For continuous variables, \emph{conditional differential entropy} $h(X|Y)$ is defined similarly:
\begin{equation}
    h(X|Y) = h(X,Y) - h(Y) = -\int \int f(x,y) \log_2 f(x|y) \, dx \, dy.
\end{equation}


\emph{Mutual information} $I(X;Y)$ quantifies the amount of information that one random variable contains about another, or the reduction in uncertainty of $X$ due to knowledge of $Y$:
\begin{equation}
    I(X;Y) = H(X) + H(Y) - H(X,Y) = H(X) - H(X|Y) = H(Y) - H(Y|X).
    \label{eq:mutualinfo}
\end{equation}

Equation \ref{eq:mutualinfo} is symmetric ($I(X;Y) = I(Y;X)$) and non-negative ($I(X;Y) \geq 0$), with $I(X;Y) = 0$ if and only if $X$ and $Y$ are independent. The maximum value occurs when one variable is a deterministic function of the other, in which case $I(X;Y) = H(X) = H(Y)$ if they have the same entropy. Mutual information can also be expressed in terms of Kullback-Leibler (KL) divergence, a measure of how one probability distribution differs from another:
\begin{equation}
    I(X;Y) = D_{KL}(p(x,y) || p(x)p(y)),
\end{equation}
where $D_{KL}(p||q) = \sum p(x) \log_2 (p(x)/q(x))$ for discrete distributions (or integral for continuous). This interpretation views mutual information as the divergence between the joint distribution and the product of marginals, quantifying dependence. For multiple variables, \emph{conditional mutual information} $I(X;Z|Y)$ is defined as:
\begin{equation}
    I(X;Z|Y) = H(X|Y) - H(X|Y,Z) = H(X,Y) + H(Y,Z) - H(X,Y,Z) - H(Y),
\end{equation}
measuring the information $Z$ provides about $X$ given $Y$. This is key in concepts like Markov chains, where $I(X;Z|Y) = 0$ if $Y$ separates $X$ and $Z$ in the chain. Similarly, in continuous settings, mutual information $I(X;Y)$ is:
\begin{equation}
    I(X;Y) = h(X) + h(Y) - h(X,Y).
\end{equation}
%which is invariant under homeomorphic transformations of the variables, making it a robust measure of dependence beyond linear correlations (e.g., unlike Pearson correlation).

Mutual information has numerous applications, such as feature selection in machine learning (where high $I(X;Y)$ indicates informative features $X$ for target $Y$), channel capacity in communication theory ($C = \max_{p(x)} I(X;Y)$ for a channel with input $X$ and output $Y$), and rate-distortion theory. For example, in a binary symmetric channel with crossover probability $p$, the mutual information is $I(X;Y) = 1 - H_b(p)$, where $H_b(p) = -p\log_2 p - (1-p)\log_2 (1-p)$ is the binary entropy function. In distributed coding contexts, mutual information bounds the achievable rates, as seen in the Slepian-Wolf and Wyner-Ziv theorems, where correlations are captured via $I(X;Y)$.

\emph{Entropy coding} refers to techniques that compress data by assigning shorter codes to more probable symbols, approaching the entropy bound. Classic methods include:

\begin{itemize}
    \item \textbf{Huffman coding}: A prefix-free code that builds a binary tree based on symbol probabilities, assigning shorter paths (codes) to frequent symbols. It achieves near-optimal compression for independent symbols but is suboptimal for correlated data.
    \item \textbf{Arithmetic coding}: A more efficient method that encodes an entire message as a single fractional number in $[0,1)$, dynamically partitioning the interval based on cumulative probabilities. It can approach the entropy limit more closely than Huffman coding, especially for sources with memory.
\end{itemize}
%These methods are lossless, meaning the original data can be perfectly reconstructed.

\subsubsection{Distributed Source Coding}
Based on the original, old paper, \cite{SlepianWolf}, shows that if you have 2 correlated sources, $X$ and $Y$, you can compress them separately while reaching the same rate as if you compressed them together. This is known as the Slepian-Wolf theorem. Wyner-Ziv extends this to lossy compression, allowing for some distortion in the reconstruction



% sub subsection - Conventional vs. Learned DSC Methods
its faster. much faster. the paper \cite{OzyilkanRecoverBinning} explains more about the speed issue in the first pages.



\subsection{Distributed Learning Primer}


% • Communication bottlenecks and gradient compression


% • Federated learning (FL), FedSGD/FedAvg, ...
